% ============================================================
%  §3  Formal Mapping
% ============================================================

\section{Formal Mapping}
\label{sec:formal}

This section establishes the correspondence between VED$\times$IFGT
variables and conceptual phenomena.
All variables follow the Master Variable Table v2.1.

\subsection{The Background Field}

The information density field $I(x,\tau) = f(1-C(x,\tau))$ provides
the continuous substrate on which conceptual structure forms.
In this paper, $I$ is not a new cognitive primitive but a projected
realization of the IFGT information-density variable in a cognitive
observation window.
In the cognitive domain, $I$ represents the density of unresolved
informational tension at position $x$ and time $\tau$: regions where
difference chains persist as quasi-closures, neither fully stabilized
nor fully dissipated.

High $I$ corresponds to conceptually active regions; low $I$ corresponds
to regions where structure has largely closed.
The strict positivity $I(x,\tau)>0$ excludes complete closure as an
ordinary state of the conceptual landscape.
It does not by itself guarantee a nonzero local gradient at every point;
gradients arise from residual asymmetry, nonuniform history, and boundary
structure within the generated field.

The information potential
\begin{equation}
  \Phi(x,\tau) = (K * I)(x,\tau)
  \label{eq:Phi_primary}
\end{equation}
is the primary geometric object.
It encodes the accumulated bias produced by prior information flow,
expressed as a scalar field whose gradient determines the direction
of least resistance for subsequent trajectories.

\subsection{Variable Correspondence Table}

\begin{center}
\renewcommand{\arraystretch}{1.5}
\begin{tabular}{p{0.22\textwidth}p{0.48\textwidth}p{0.22\textwidth}}
\toprule
\textbf{VED$\times$IFGT variable} &
\textbf{Conceptual domain interpretation} &
\textbf{Formal role} \\
\midrule
$C(x,\tau)$ &
Degree of conceptual fixation at $x$ &
Background field; $C \in [0,\Cmax)$ \\
$I(x,\tau) = f(1-C)$ &
Density of unresolved conceptual tension &
Central field; $I > 0$ always \\
$\Phi = K*I$ &
The conceptual landscape itself &
Potential geometry \\
$\nabla\Phi$ &
Direction of least cognitive resistance &
Constraint force on $\JI$ \\
$\JI = -D\nabla I + \mu F_I$ &
Thought as flow &
Trajectory field \\
$\sigma = \sigma_d - \sigma_p - \sigma_b$ &
Rate of landscape deformation &
Source term \\
$\Delta$ &
Novel experience, perception, or linguistic input &
Driving source \\
$\delta I$ &
Unresolved conceptual residue &
Trace that constitutes $I$ \\
$\tau$ &
Experiential time of conceptual development &
Irreversible ordering \\
$\ell_i = \mathcal{Q}[I,\Phi,C;\cdot]$ &
Reusable linguistic token (natural language word, symbol, emblem,
gesture, diagram, formula, or any discrete unit capable of acting on the
field as a localized perturbation) &
Discrete operator on $\Phi$ \\
$\Phistar$ &
Mature, sedimented conceptual landscape &
Long-time attractor of $\Phi$ \\
$\Cmax$ &
Ideal complete-closure boundary &
Boundary representation, not an ordinary state \\
\bottomrule
\end{tabular}
\end{center}

\begin{remark}[Notational remark]
The operator $\mathcal{Q}$ is called a quantization operator because its
action --- extracting discrete, re-addressable units from a continuous
field configuration --- is structurally isomorphic to quantization in
the physical sense: mapping continuous variables to a discrete set of
stable values.
The term ``quantum'' is not used as a theoretical claim within CT, but
the analogy is preserved in the name of this operator.
\end{remark}

\subsection{What the Mapping Does Not Claim}

The mapping established above does not introduce semantic content,
intentionality, or psychological states into the formalism.
In particular:
\begin{itemize}
  \item $\Phi$ is not a meaning field. It is a dynamical constraint
        geometry that determines which transitions are cheap and which
        are costly.
  \item $\JI$ is not thought content. It is a flow direction; the
        content that appears in experience is a downstream phenomenon.
  \item $\ell_i$ is not a meaning unit. It is a discrete re-addressable
        unit that couples locally to $\nabla\Phi$.
\end{itemize}
These exclusions are methodological necessities, not claims about the
ultimate nature of meaning or experience.
